% VERSION WITHOUT SOLUTIONS
\documentclass{jn-pset}
% Uncomment to get VERSION WITH SOLUTIONS and also do
%   \setboolean{versionwithsolutions}{true}
% a few lines down
%  cumentclass[solutions]{jn-pset}

\usepackage{ifthen}
\newboolean{versionwithsolutions}
% Uncomment to insert text that should only be there in version WITH solutions
\setboolean{versionwithsolutions}{true}
% Uncomment to insert text that should only be there in version WITHOUT solutions
\setboolean{versionwithsolutions}{false}



%%%%
%%%% QUICK INSTRUCTIONS FOR FORMATTING OF PROBLEMS
%%%%
%    
% Code a problem by
%    \begin{problem}
%    \end{problem}
%
% Code a subproblem inside a problem by (note the percent signs!)
%    \begin{subproblem}%
%        \label{problem:labelhere}%
%        Text here
%    \end{subproblem}
%
% Get a small vertical space by issuing the command
%    \smallskip
%
% For instance, to give a hint to a problem (after having completed
% the problem statement) code in the following way:
%
%    \smallskip
%     \noindent
%     \emph{Hint:}
%     Consider the following super-useful hint for this particular problem...
%

% PROBLEM-SET-SPECIFIC MACROS (UPDATE FOR EACH PROBLEM SET)
% *** NOT USED FOR EXAM ***
%
% \newcommand{\psetno}{5}
%%% The formatting macro \duedate is not used for the exam
%   \newcommand{\duedate}{Wednesday April 10 at 13:00 CET}
% Dummy formatting for draft versions
%   \newcommand{\thresholdA}{$\langle A \rangle$~points\xspace}
%   \newcommand{\thresholdB}{$\langle B \rangle$~points\xspace}
%   \newcommand{\thresholdC}{$\langle C \rangle$~points\xspace}
%   \newcommand{\thresholdD}{$\langle D \rangle$~points\xspace}
%   \newcommand{\thresholdE}{$\langle E \rangle$~points\xspace}
% Add actual grading thresholds here with \renewcommand 
%   \renewcommand{\thresholdA}{$200$~points\xspace}
%   \renewcommand{\thresholdB}{$170$~points\xspace}
%   \renewcommand{\thresholdC}{$140$~points\xspace}
%   \renewcommand{\thresholdD}{$110$~points\xspace}
%   \renewcommand{\thresholdE}{$80$~points\xspace}

% Threshold guaranteed to pass the exam
\newcommand{\thresholdforpass}{$220$~points\xspace}


% Compact lists
\usepackage{paralist}

% COURSE-SPECIFIC MACROS

\input{../ProblemSets/course-macros.tex}

\ifthenelse{\boolean{versionwithsolutions}}
{\newcommand{\examheading}{\mbox{Exam April 8, 2026} \\
    Problems and Solutions}}
{\newcommand{\examheading}{\mbox{Exam April 8, 2026}}}


% For getting watermark "DRAFT" across all pages (for instance, 
% when posting preliminary version of problem set)
%   \usepackage{draftwatermark}
%   % \SetWatermarkFontSize{20 cm}
%   \SetWatermarkScale{5}


%%%
%%% TITLE
%%%

\author{\courseinstructor}
\course{\coursenamelong{}}
\semester{\courseperiod}
\title{\coursenameshort: 
  \examheading}



\begin{document}

\maketitle

\ifthenelse{\boolean{versionwithsolutions}}
{
\noindent
\emph{%
  Please note that the main purpose of this document is to explain
  what the correct solutions are and how to arrive at them.  Thus,
  while the text below is certainly intended to provide good examples
  of how to solve problems and reason about solutions, these examples
  do not necessarily specify exactly how the handed-in exams were
  expected to look like.
%
  This is especially so since for many problems there are more than
  one correct way of solving them.  
%   
%     In general, the solutions provided below tend to be much more
%     detailed than what would be expected for the exam submissions.
%   
  As communicated during the course, the course notes
  published on Absalon, which contain many problems that we worked out
  in class, are probably the best indicator of what level of detail is
  expected from the exam solutions.
}


}
{\input{IDMA_Exam_GeneralInstructions.tex}}

%%%
%%%   MATHEMATICAL INDUCTION PROBLEM
%%%
\begin{problem}
  (60 p)
%     Provide formal proofs of the following claims using proof
  Solve the following problems using
  techniques that we have learned during the course.
  Please make sure to provide not just intuition but also formal proofs.

  \begin{subproblem}%
    \ifthenelse{\boolean{versionwithsolutions}}
    {(30 p)}
    {\ignorespaces}
    Prove that if
    $x_1, x_2, \ldots, x_n$
    are non-negative real numbers, then
    it always holds that
%       the inequality
    \begin{equation*}
      (1 + x_1) (1 + x_2) \cdots  (1 + x_n)
      \geq
      1 + x_1 + x_2 + \cdots + x_n
      \eqperiod
    \end{equation*}
%       always holds.
  \end{subproblem}

\begin{solution}
We prove this inequality by induction over~$n$.

\begin{description}
\item[\emph{Base case ($n=1$):}]
  In this case we have that equality
  $
  1 + x_1
  =
  1 + x_1
  $
  holds between the left-hand side and the right-hand side.
  
\item[\emph{Induction step:}]
\end{description}
The claim now follows by the induction principle.  
\end{solution}


\begin{subproblem}%
  \ifthenelse{\boolean{versionwithsolutions}}
  {(30 p)}
  {\ignorespaces}
  Prove that every integer
  $n \geq 12$
  can be written as
  $4a + 5b$ for $a$ and $b$ non-negative integers.
\end{subproblem}

\begin{solution}
%     This is very similar to a problem we did in class, 
%     and we will solve this one in the same way by arguing by
  We prove this by 
  induction over~$n$.
  \begin{description}
  \item[\emph{Base cases:}]
    Here we need to consider several base cases:
    
    
  \item[\emph{Induction step:}]
  \end{description}
  The claim follows by the induction principle.
\end{solution}


\end{problem}


%%%
%%% COMBINATORICS
%%%

\begin{problem}
  (60 p)
  As Jakob was preparing this year's instalment of the IDMA course in
  early $2026$, he realized that he had been teaching discrete
  mathematics and algorithms in Copenhagen for $6$~years.  For some
  reason that is not so easy to explain, this made Jakob wonder what the
  largest positive integer~$n$ is such that
  $26! / 6^n$
  is also an integer.
  Could you help Jakob figure this out?  Jakob does not
  really trust calculators, so he would very much like to see a clean and simple
  solution by hand that just uses what we learned during the IDMA course.
\end{problem}

\begin{solution}
  SOLUTION HERE
\end{solution}

%%% 
%%% PSEUDOCODE PROBLEM LINKED LISTS
%%%

\begin{problem}
  (80 p)
  In this problem we consider linked lists.
  The records in our list datatype have four attributes:
  \begin{itemize}
  \item
    \verb+data+: the data stored in the record (which we will not care about, as usual);
    
  \item
    \verb+key+: used for sorting the records in the list;
    
  \item
    \verb+prev+: link to the previous element in the list;
    
  \item
    \verb+next+: link to the next element in the list.
    
  \end{itemize}
%%%
  The lists are sorted in increasing order based on the key values in
  \verb+key+.
  We assume that distinct records in our list have distinct
  key values, so that every value identifies a unique record in a list.
  To make the coding easier, every list is initialized with two
  dummy records  
  \verb+head+
  and
  \verb+tail+
  that do not contain any data but are
  such that
  \begin{itemize}
  \item
    \verb+head.key+ $=$ \verb+-INFINITY+,
    \verb+head.prev+ $=$ \verb+NULL+,
    and
    \verb+head.next+ $=$ \verb+tail+;
    
  \item
    \verb+tail.key+ $=$ \verb|+INFINITY|,
    \verb+tail.prev+ $=$ \verb+head+,
    and
    \verb+tail.next+ $=$ \verb+NULL+.
  \end{itemize}
  %%% 
  This ensures that \verb+head+ will always be the first element
  in the list since
  \verb|-INFINITY|$= - \infty$ is smaller than any key value
  and that 
  \verb+tail+ will always be last  since
  \verb|+INFINITY|$= +\infty$ is larger than any key value.

  Below you find four pseudocode methods for operating on the list.
  Please identify which method:
  \begin{enumerate}[(a)]
  \item
    \textbf{Inserts} an element
    identified by a key value
    in the list.
    
  \item
    \textbf{Removes} an element
    identified by a key value
    from the list.

  \item
    \textbf{Looks up and returns}
    an element
    identified by a key value
    from the list.

  \end{enumerate}
%%%
  Note that there is one spoiler piece of code that does none of the
  above.
  All four pieces of code start with exactly the same while loop,
  and all of them have a return statement regardless of whether
  this is needed or not (since it would be too easy to
  distinguish the methods otherwise).

  You receive 10 points for each correct match and 10 more
  points for a brief but convincing explanation why the piece of code
  does what you claim it does (or, for the spoiler piece of code, why it
  does something else and why this does not make any sense).

%%% REMOVAL
\begin{verbatim}
METHOD1 (record) 
    current := head
    while record.key > current.key {
        current := current.next
    }
    if record.key == current.key {
        previous      := current.prev
        current       := current.next
        previous.next := current
        current.prev  := previous
    }
    return record
\end{verbatim}

%%% LOOK-UP
\begin{verbatim}
METHOD2 (record) 
    current := head
    while record.key > current.key {
        current := current.next
    }
    if record.key != current.key {
        current := NULL
    }
    return current
\end{verbatim}


%%% SPOILER
\begin{verbatim}
METHOD3 (record) 
    current := head
    while record.key > current.key {
        current := current.next
    }
    if record.key >= current.key {
        previous      := current.prev
        current.prev  := previous.next
        previous      := previous.next
        current.next  := previous
    }
    return current
\end{verbatim}

%%% INSERTION 
\begin{verbatim}
METHOD4 (record) 
    current := head
    while record.key > current.key {
        current := current.next
    }
    if record.key < current.key {
        previous      := current.prev
        previous.next := record
        record.prev   := previous
        record.next   := current
        current.prev  := record
    }
    return record
\end{verbatim}



\begin{solution}
  SOLUTION HERE
\end{solution}

\end{problem}


%%%
%%% PROPOSITIONAL LOGIC
%%%

\begin{problem}%
  \label{problem:logic}%
  (60 p)  
  For each of the propositional logic formulas below,
  determine whether it is a tautology or not.
  If the formula is not a tautology, show how to
  add a single connective
  $\land$, $\lor$, or $\lnot$
  to make it into a tautology.
  Please make sure to justify your answers
  (e.g., by presenting truth tables, or by using rules for rewriting logic formulas
  that we have learned in class).

%%% TAUTOLOGY  
  \begin{subproblem}%
    \label{problem:formula-1}%
    \ifthenelse{\boolean{versionwithsolutions}}
    {(30 p)}
    {\ignorespaces}
    $
    \bigl(
    ( p \limpl q ) \limpl r 
    \bigr)
    \lor
    \bigl(
    (\lnot p \land \lnot r)
    \lor
    (q \land \lnot r)
    \bigr)
    $
  \end{subproblem}

\begin{solution}
  This formula is a tautology, i.e., it is always true.

  *** UPDATE ***
  
  To see this, note first that we know that
  \begin{equation}
    \label{eq:impl}
    p \limpl q
    \ \equiv \
    \lnot p \lor q
  \end{equation}
  since the only way the implication
  $p \limpl q$
  can be false is that $p$ is true and $q$ is false,
  and this is one of the basic equivalences we have learned in the course.
  Furthermore, from De Morgan's laws it is easy to derive that
  \begin{equation}
    \label{eq:de-morgan}
    \lnot ( a \lor b \lor c)
    \ \equiv \
    \lnot a \land \lnot b \land \lnot c
  \end{equation}
  (i.e., the only way a disjunction can be false is that all its
  disjuncts are false).
  
  By combining
  \eqref{eq:impl}
  and~\eqref{eq:de-morgan},
  and using that conjunction is commutative, we see that the premise
  $
  \lnot \bigl(
  ( p \limpl q) \lor r
  \bigr)
  $
  is in fact equivalent to the conclusion
  $
  (\lnot q \land \lnot r) \land p
  $,
  and so the implication
  in Problem~\ref{problem:formula-1}
  will always evaluate to true.
\end{solution}

  \begin{subproblem}
    \ifthenelse{\boolean{versionwithsolutions}}
    {(30 p)}
    {\ignorespaces}
    $
    \bigl(
    p \limpl (q \lor r)
    \bigr)
    \lequiv
    \bigl(
    (p \lor q)
%%% This should be    
%       (\lnot p \lor q)
    \lor
    (p \land r)
    \bigr)
    $    
  \end{subproblem}

\begin{solution}
  This is not a tautology.

  *** UPDATE ***
  If we set $p$ and~$q$ to true
  but $r$ to false, we get that
  $
  (p \land q) \limpl r
  $
  evaluates to  false but
  $
  (q \lor r) \lor \lnot p
%%% This should be    
%       (\lnot q \lor r) \lor \lnot p
  $
  evaluates to true, and so the whole formula is false.

  If we change $q$ on the right-hand side to $\lnot q$,
  we obtain the formula
  \begin{equation}
    \label{eq:patched}
    \bigl(
    (p \land q) \limpl r
    \bigr)
    \lequiv
    \bigl(
    (\lnot q \lor r) \lor \lnot p
    \bigr)
    \ ,
  \end{equation}    
  which is a tautology. To see this, we can rewrite the left-hand side
  of~\eqref{eq:patched}
  as
  \begin{subequations}
  \begin{align}
    (p \land q) \limpl r
    \ &\equiv \
        \lnot (p \land q) \lor r
    \\
    \ &\equiv \
        (\lnot p \lor \lnot q) \lor r
    \\
    \ &\equiv \
        (\lnot q \lor r) \lor \lnot p
  \end{align}
  \end{subequations}
  by using
  the property~\eqref{eq:impl} of implication,
  De Morgan's laws,
  and commutativity of disjunction.
  This show that the left-hand and right-hand sides of
  of the formula~\eqref{eq:patched} are equivalent,
  and so the formula will always evaluate to true.
\end{solution}

\end{problem}



%%% 
%%% HEAPSORT
%%%

\providecommand{\buildmaxheap}{\textsc{Build-Max-Heap}\xspace}
\providecommand{\maxheapify}{\textsc{Max-Heapify}\xspace}

\begin{problem}
  (70 p)
  In this problem we wish to study heaps and how they can be used for sorting.

\begin{subproblem}%
  \label{problem:heap-sort}%
  \ifthenelse{\boolean{versionwithsolutions}}
  {(50 p)}
  {\ignorespaces}
  Run heapsort by hand to sort the array
%%%
  $
  A =
  [ 11, 3, 19, 5, 17, 2, 13, 7 ]
  $
%%%
  in increasing order
  and give a detailed explanation of the first few steps of the
  algorithm.
  Specifically,
  show how the
  % original
  heap is built and how the first two numbers
%     in the array
  are sorted into their correct final positions.
  Make sure to explain what calls to the heap are made from the
  sorting algorithm and what comparisons are made, 
  and illustrate how the heap and array have changed at the
  end of every ``heapify'' or ``bubble'' call (after all recursive
  subcalls have been taken  care of).

  \smallskip
  \noindent
  \emph{Remark:}
  Please note that
%     you do \emph{not} have to
  there is no need to
  present the entire execution of the sorting algorithm.
%     , since this might be quite time-consuming.
  Just showing the initial steps
  as described above is fully sufficient.
\end{subproblem}

\begin{solution}
  ***  UPDATE ***
  We illustrate in
  \reffig{fig:sorting-heaps}
  how the heap used for sorting the array is built
  and then modified as the first two numbers are extracted.
%
  \begin{enumerate}
  \item 
    The input array is shown in a heap structure in
    \reffig{fig:sorting-heap-1}.
    
    \input{Figures/figHeapsSort.tex}  
    
  \item 
    We need to run \buildmaxheap, which builds max-heaps bottom-up
    using the \maxheapify method. The $6$th through $10$th nodes
    vacuously constitute max-heaps of size~$1$, so the first time
    \maxheapify is run is on the $5$th node containing~$3$.
    Since the child~$10$ is larger, the elements
    $3$ and~$10$ swap places as shown in
    \reffig{fig:sorting-heap-2}.
    When \maxheapify is called for the subheap rooted at the $10$th
    node, it returns immediately, since that subheap has size~$1$
    and so is a correct max-heap vacuously. 

  \item 
    \maxheapify is then run on the $4$th node. The largest child~$9$
    is larger than the number~$7$ stored in the $4$th node, so they
    trade places as shown in
    \reffig{fig:sorting-heap-3}. Again, the second recursive
    \maxheapify call returns immediately.

  \item 
    For the $3$rd node, the number~$4$ is smaller than the largest
    child~$8$, so the two numbers are swapped and the recursive
    \maxheapify call returns immediately, resulting in the heap in
    \reffig{fig:sorting-heap-4}.

  \item 
    For the $2$nd node, the number~$6$ is smaller than the largest
    child~$10$, so they exchange positions with $6$ being shifted to
    the $5$th node.
    For the recursive \maxheapify call on the $5$th node, the
    number~$6$ is larger than the only child~$3$ and so the subheap is
    in order. After this the heap looks as in
    \reffig{fig:sorting-heap-5}.

  \item 
    In the final \maxheapify call of \buildmaxheap, the number~$5$ at the
    root is compared to both children of the root. The largest
    child~$10$  is moved to the root and$5$~is shifted down to the
    $2$nd~node. 
    In the first recursive \maxheapify call, $5$ is compared to the
    children $9$ and~$6$, and the largest child~$9$ is moved up and
    $5$~is shifted down to the $4$th~node.
    In the next \maxheapify call for the $4$th~node,
    $5$ is swapped with the largest child~$7$. The final recursive call
    for the $8$th node terminates since the subheap has size~$1$.
    The max-heap resulting from the \buildmaxheap call is shown in
    \reffig{fig:sorting-heap-6}.

  \item
    The heapsort algorithm now repeatedly removes the root, places
    the last element in the heap in the root position, and runs
    \maxheapify.
    Right after $10$ has been removed and $3$ has been put in the root
    position, we have the heap structure in    
    \reffig{fig:sorting-heap-7}.
    A sequence of recursive \maxheapify calls compare
    $3$  first  to the largest of the two children $9$ and~$8$ of the
    root,
    moving up~$9$ and moving down~$3$ to the $2$nd node,
    then to the largest of the two children $7$ and~$6$,
    moving up~$7$ and moving~$3$ further down to the $4$th~node, 
    and finally to the largest of the two children $5$ and~$1$, 
    swapping $3$ and~$5$.
    This yields the max-heap on nine elements in
    \reffig{fig:sorting-heap-8}.

  \item
    Right after $9$ has been removed and $1$ has been put in the root
    position, we have the heap structure in    
    \reffig{fig:sorting-heap-9}.
    The recursive \maxheapify calls compare
    $1$  first  to the largest of the two children $7$ and~$8$ of the root,
    changing places of $1$ and $8$, 
    and then to the largest of the two children $4$ and~$2$, 
    changing places of $1$ and $4$. 
    The resulting  max-heap on eight elements is depicted in
    \reffig{fig:sorting-heap-8}.

  \end{enumerate}
%
  This completes the description of how the first two numbers $10$ and
  $9$ are extracted from the heap during heapsort, which is what we
  were asked to explain and illustrate.
\end{solution}


\begin{subproblem}
  \ifthenelse{\boolean{versionwithsolutions}}
  {(20 p)}
  {\ignorespaces}
  Jakob thinks he has figured out a nifty way
  of finding the \emph{minimal} element in a max-heap by calling the
  following recursive algorithm with index $i = 1$:
\begin{verbatim}
MAX-HEAP-FIND-MIN (A, i)
    if (2 * i > A.heapsize) {
        smallest := i
    }
    else {
        smallest := 2 * i
    }
    if (2 * i + 1 <= A.heapsize AND A[2 * i + 1] < A[smallest]) {
        smallest := 2 * i + 1
    }
    if (smallest == i) {
        return (A[i])
    }
    else {
        return (MAX-HEAP-FIND-MIN (A, smallest))
    }
\end{verbatim}
%
  Is Jakob right that this will always return the smallest element
  (assuming that $A$ is a correct max-heap)?
  Please provide either a proof of correctness or a counterexample
  where the input is a max-heap but
%     Jakob's algorithm
  \verb+MAX-HEAP-FIND-MIN+
  fails to find the
  minimal element. 
\end{subproblem}

\begin{solution}
  No, this is not true.

*** UPDATE ***

  If we consider the array from our max-heap in 
  Problem~\ref{problem:heap-sort},
  it looks like
  $
  10,9,8,7,6,4,2,5,1,3
  $
  (see again \reffig{fig:sorting-heap-6}).
  If we take the array in the reverse order to get
  $
  3,1,5,2,4,6,7,8,9,10
  $,
  then we have that $1$ is the left child of the root $3$, which
  violates the min-heap property. 
\end{solution}


\end{problem}



%%%
%%% TRICKY INDUCTION PROBLEM INVOLVING FIBONACCI NUMBERS
%%%

\begin{problem}%
  \label{problem:grid-tilings}%
  (80 p)
  In this problem we wish to count the number of ways of tiling a
  $2 \times n$ grid
  (as in
  \reffig{fig:tiling-2-by-n-grid-example})  
  with
  $2 \times 1$ tiles
  so that every cell in the grid is covered by exactly one tile.
%%%  
%     (that can be rotated so that they have either
%     height~$2$ and length~$1$
%     or
%     height~$1$ and length~$2$).
%%%
%     It is not so hard to see that there are:
%     \begin{itemize}
%         
%     \item
%       $1$ tiling for a $2 \times 1$ grid
%       (\reffig{fig:tiling-2-by-1-grid});
%     \item
%       $2$ different tilings for a $2 \times 2$ grid
%       (\reftwofigs{fig:tiling-2-by-2-grid-a}{fig:tiling-2-by-2-grid-b});
%       
%     \item
%       $3$ different tilings for a $2 \times 3$ grid
%       (\refthreefigs
%       {fig:tiling-2-by-3-grid-a}
%       {fig:tiling-2-by-3-grid-b}
%       {fig:tiling-2-by-3-grid-c}).
%       
%     \end{itemize}
%   
%%%
%     It is not so hard to see that there is
  There is
  $1$ such tiling for a $2 \times 1$ grid
  (\reffig{fig:tiling-2-by-1-grid});
  exactly $2$~different tilings for a $2 \times 2$ grid
  (\reftwofigs{fig:tiling-2-by-2-grid-a}{fig:tiling-2-by-2-grid-b});
  and
  exactly
  $3$~different tilings for a $2 \times 3$ grid
  (\refthreefigs
  {fig:tiling-2-by-3-grid-a}
  {fig:tiling-2-by-3-grid-b}
  {fig:tiling-2-by-3-grid-c}).

  
\begin{figure}[tp]

  \begin{subfigure}[b]{1\textwidth}
  \centering
    \includegraphics{Figures/tiling-2-by-n-grid.4}%      
    \caption{Example $2 \times 20$  grid.}
    \label{fig:tiling-2-by-n-grid-example}    
  \end{subfigure}

  \begin{subfigure}[b]{0.3\textwidth}
    \centering
    \includegraphics{Figures/tiling-2-by-n-grid.5}%          
    \caption{%
      Only
      tiling of
      $2 \times 1$ grid.
    }	    
    \label{fig:tiling-2-by-1-grid}
  \end{subfigure}
  \hfill
  \begin{subfigure}[b]{0.3\textwidth}
    \centering
    \includegraphics{Figures/tiling-2-by-n-grid.6}%          
    \caption{%
      First
      tiling of
      $2 \times 2$ grid.
    }	    
    \label{fig:tiling-2-by-2-grid-a}
  \end{subfigure}
  \hfill
  \begin{subfigure}[b]{0.3\textwidth}
    \centering
    \includegraphics{Figures/tiling-2-by-n-grid.7}%          
    \caption{%
      Second      
      tiling of
      $2 \times 2$ grid.
    }	    
    \label{fig:tiling-2-by-2-grid-b}
  \end{subfigure}
  

  \begin{subfigure}[b]{0.3\textwidth}
    \centering
    \includegraphics{Figures/tiling-2-by-n-grid.8}%          
    \caption{%
      First
      tiling of
      $2 \times 3$ grid.
    }	    
    \label{fig:tiling-2-by-3-grid-a}
  \end{subfigure}
  \hfill
  \begin{subfigure}[b]{0.3\textwidth}
    \centering
    \includegraphics{Figures/tiling-2-by-n-grid.10}%          
    \caption{%
      Second
      tiling of
      $2 \times 3$ grid.
    }	    
    \label{fig:tiling-2-by-3-grid-b}
  \end{subfigure}
  \hfill
  \begin{subfigure}[b]{0.3\textwidth}
    \centering
    \includegraphics{Figures/tiling-2-by-n-grid.9}%          
    \caption{%
      Third
      tiling of
      $2 \times 3$ grid.
    }	    
    \label{fig:tiling-2-by-3-grid-c}
  \end{subfigure}
  
  \caption{Tilings of
    $2 \times n$ grids in
    Problem~\ref{problem:grid-tilings}.}
  \label{fig:tiling-2-by-n-grid}
\end{figure}


  
  \begin{subproblem}
      \ifthenelse{\boolean{versionwithsolutions}}
    {(20 p)}
    {\ignorespaces}
    How many different tilings are there for a $2 \times 4$ grid?
    Make sure to explain how you compute this number so that we can
    follow your reasoning.
  \end{subproblem}

  \begin{subproblem}
      \ifthenelse{\boolean{versionwithsolutions}}
    {(60 p)}
    {\ignorespaces}
    How many different tilings are there for a $2 \times n$ grid
    for a positive integer~$n$?

    \smallskip
    \noindent
    \emph{Hint:}
%       We want to find
    Find a nice formula for this number
    in terms of concepts we learned in class.
  \end{subproblem}


\begin{solution}


  \begin{description}
  \item[\emph{Base case ($n=2$):}]

    
  \item[\emph{Induction step:}]

  \end{description}
%%%  
  The claim in the problem statement now follows by the induction principle.
\end{solution}


\end{problem}



%%%
%%% STRONGLY CONNECTED COMPONENTS (AND RELATIONS)
%%%

\begin{problem}%
  \label{problem:scc}%
  (80 p)  
%%%
%%% a: 10 p
%%% b: 10 p
%%% c: 20 p
%%% d: 40 p
  This problems is about relations in general and
  graphs in particular.
%   
%     We will focus on the graph in
%     \reffig{fig:scc}.
%   
%     \input{Figures/fig-scc-26.tex}    
%   
  
  \begin{subproblem}
    \ifthenelse{\boolean{versionwithsolutions}}
    {(10 p)}
    {\ignorespaces}
  \label{problem:scc-a}%
    Say that the relation $E(u,v)$ holds for two vertices $u$ and $v$
    in a directed graph if there is
%       a directed
    an
    edge from $u$ to~$v$.   
    Present the matrix representation
%       of this relation
    of~$E$
    for the graph in
    \reffig{fig:scc}.
%       (with rows and columns sorted in lexicographic order of the vertices).
  \end{subproblem}

  \input{Figures/fig-scc-26.tex}    
  
\begin{solution}
  \begin{equation}
    \label{eq:adjacency-matrix}
    M_E =
    \begin{pmatrix}
       0 & 1 & 1 & 1 & 0 & 0 \\
       0 & 0 & 1 & 1 & 0 & 0 \\
       0 & 0 & 0 & 0 & 0 & 1 \\
       0 & 0 & 0 & 0 & 1 & 0 \\
       0 & 0 & 0 & 1 & 0 & 0 \\
       0 & 1 & 1 & 0 & 0 & 0 
    \end{pmatrix}
  \end{equation}
\end{solution}

  \begin{subproblem}
    \ifthenelse{\boolean{versionwithsolutions}}
    {(10 p)}
    {\ignorespaces}
    \label{problem:scc-b}%
    Take the reflexive closure $R$ of the relation~$E$ in
    Problem~\ref{problem:scc-a}
    and 
    present the matrix representation of this relation.
  \end{subproblem}

\begin{solution}
  \begin{equation}
    \label{eq:adjacency-matrix-reflexive-closure}
    M_R =
    \begin{pmatrix}
       1 & 1 & 1 & 1 & 0 & 0 \\
       0 & 1 & 1 & 1 & 0 & 0 \\
       0 & 0 & 1 & 0 & 0 & 1 \\
       0 & 0 & 0 & 1 & 1 & 0 \\
       0 & 0 & 0 & 1 & 1 & 0 \\
       0 & 1 & 1 & 0 & 0 & 1 
    \end{pmatrix}
  \end{equation}
\end{solution}

  \begin{subproblem}
    \ifthenelse{\boolean{versionwithsolutions}}
    {(20 p)}
    {\ignorespaces}
    \label{problem:scc-c}%
%       
    Take the transitive closure $T$ of the relation~$R$ in
    Problem~\ref{problem:scc-b}
    and 
    present the matrix representation of this relation.
%
%       Present the matrix representation of
%       the transitive closure $T$ of the relation~$R$ in
%       Problem~\ref{problem:scc-b}.
  \end{subproblem}

  
\begin{solution}
  \begin{equation}
    \label{eq:adjacency-matrix-transitive-closure}
    M_T =
    \begin{pmatrix}
       1 & 1 & 1 & 1 & 1 & 1 \\
       0 & 1 & 1 & 1 & 1 & 1 \\
       0 & 1 & 1 & 1 & 1 & 1 \\
       0 & 0 & 0 & 1 & 1 & 0 \\
       0 & 0 & 0 & 1 & 1 & 0 \\
       0 & 1 & 1 & 1 & 1 & 1 
    \end{pmatrix}
  \end{equation}
\end{solution}

  \begin{subproblem}
    \ifthenelse{\boolean{versionwithsolutions}}
    {(20 p)}
    {\ignorespaces}
    \label{problem:scc-d}%
    Now consider a new relation $S(u,v)$
    where we define 
    $S(u,v)$ to hold for two vertices $u$ and~$v$ precisely when
    $T(u,v)$ and~$T(v,u)$ both hold
    for the relation~$T$ defined as in 
    Problem~\ref{problem:scc-c}.
        
    This gives a particular type of relation that we discussed in class.
    What is this type of relation called?
    Explain, briefly but clearly, why we are guaranteed to always get
    this type of relation when performing the steps described in
    Problems~%
    \ref{problem:scc-b},
    \ref{problem:scc-c},
    and~%
    \ref{problem:scc-d}.

    Also, this relation~$S$ groups the vertices of the graph in a
    particular way. What is the technical term that we use for such
    groups of vertices in a graph?    

    \smallskip
    \noindent
    \emph{Remark:}
    Note that you do \emph{not} have to present any matrix
    representation for this relation~$S$.
%       We are just looking for answers to the questions above.
    
  \end{subproblem}

  \begin{solution}
    Let us give the matrix to be helpful
  \begin{equation}
    \label{eq:adjacency-matrix-transitive-closure}
    M_S =
    \begin{pmatrix}
       1 & 0 & 0 & 0 & 0 & 0 \\
       0 & 1 & 1 & 0 & 0 & 1 \\
       0 & 1 & 1 & 0 & 0 & 1 \\
       0 & 0 & 0 & 1 & 1 & 0 \\
       0 & 0 & 0 & 1 & 1 & 0 \\
       0 & 1 & 1 & 0 & 0 & 1 
    \end{pmatrix}
  \end{equation}
\end{solution}
  
%%%%
%%%% WILL IT JUST BE TOO MUCH TO ASK THEM TO RUN SCC?
%%%%
%   
%   - In fact, we have even seen an algorithm for computing this
%   relation. Run this algorithm on our small example graph, in the same
%   way as we did in class,
%   
\end{problem}



%%%
%%% TRICKY COMBINATORICS
%%%

\begin{problem}%
  \label{problem:multisets}%
  (100 p)  
%%%
  Suppose that we have a universe of $n$ elements and are interested
  in counting the number of multisets of size $r$.
  Jakob claims to have figured out that the number of such multisets is~%
  $\binom{n + r - 1}{n-1}$.
  His idea is to prove this by
  using the following algorithm to translate
  any given size-$r$ multiset~$M$ of~$[n]$
  to a standard subset~$S$ of $[n + r - 1]$ of size $n-1$:
%     (which assumes that the elements in the multiset are sorted
%     in increasing order): 
\begin{verbatim}
MULTISET2SET (M)
    elem := 1
    S    := empty-set
    for i := 1 upto n - 1 {
        c := number of copies of element i in M
        for j := 1 upto c {
            elem := elem + 1
        }
        S := S union {elem}
        elem := elem + 1
    }
    return S
\end{verbatim}

  \begin{subproblem}%
    \label{problem:multiset-1}%
    \ifthenelse{\boolean{versionwithsolutions}}
    {(20 p)}
    {\ignorespaces}
    Suppose that we have $n = 5$ elements in the universe and are interested in
    multisets of \mbox{size $r = 6$}.
%   
%       Describe what subset $S$ the multiset $M = [1, 1, 2, 5, 5, 5]$ is
%       translated to by Jakob's algorithm.
%   
    Describe what subset $S$ is returned when the algorithm
    \verb+MULTISET2SET+
    above is run with the multiset $M = [1, 1, 2, 5, 5, 5]$ as input.
    Make sure to explain, briefly but clearly,
    how this set is computed.
  \end{subproblem}

\begin{solution}
  SOLUTION HERE
\end{solution}

  \begin{subproblem}%
    \label{problem:multiset-2}%
    \ifthenelse{\boolean{versionwithsolutions}}
    {(40 p)}
    {\ignorespaces}
%%%
    Is Jakob's new algorithm
    \verb+MULTISET2SET+
    correct? Can you prove that the
    algorithm provides a
    translation that shows that the number of multisets of size $r$ from
    a universe of size $n$ is exactly $\binom{n + r - 1}{n-1}$? Or if there is an
    error in Jakob's algorithm, then what is it?

    \smallskip
    \noindent
    \emph{Hint:}
    A proof of correctness would need to go through essentially the
    same kind of steps that we did in class. If Jakob's method is in
    fact incorrect, then just one small counterexample is sufficient.
  \end{subproblem}

\begin{solution}
  SOLUTION HERE
\end{solution}


  \begin{subproblem}%
    \label{problem:multiset-3}%
    \ifthenelse{\boolean{versionwithsolutions}}
    {(40 p)}
    {\ignorespaces}
    Consider the original translation from multisets of size $r$ from a
    universe of size $n$ to subsets of size $r$ from a universe of
    size
    $n + r - 1$ that we used in class, and that we proved to be correct.
%%%
    For fixed positive integers $n$ and~$r$, let us say that the Boolean predicate
    $T_{n,r}(M,S)$
    is true if the multiset $M$ can be translated to the standard
    subset $S$ by the method we studied in class, and that
    $T_{n,r}(M,S)$
    is false otherwise.
    
%       What properties does this predicate $T_{n,r}$ have?
    Please indicate
    for all alternatives below whether $T_{n,r}$ satisfies this property or
    not. For each property, you get points for the correct
    answer,
%       whether $T_{n,r}$ has the property or not,
    a deduction for
    a wrong answer, and no change in your score if you do not
    indicate an answer for a specific property.
    \begin{enumerate}  %  [(i)]
    \item
      $T_{n,r}$ describes a relation.
      
    \item
      $T_{n,r}$ describes a function.
      
    \item
      $T_{n,r}$ is surjective.
      
    \item
      $T_{n,r}$ is reflexive.
  
    \item
      $T_{n,r}$ is a partial order.
      
    \item
      $T_{n,r}$ is injective.
      
    \item
      $T_{n,r}$ is symmetric.
      
    \item
      $T_{n,r}$ is an equivalence relation.
      
    \end{enumerate}
    
  \end{subproblem}

\begin{solution}
  SOLUTION HERE
\end{solution}

\end{problem}



%%% 
%%% GRAPH PROBLEM: IDENTIFY OUTPUT OF DIFFERENT ALGORITHMS
%%%
%   Give 10 points per correct answer.
%   Give 20 extra points for a overview of how and why the algorithms process the vertices in this order.
%   

\begin{problem}%
  \label{problem:graph-algorithm-outputs}%
  (120 p)
  Consider the graph in   
  \reffig{fig:graph-algorithm-outputs}
  and the following ordered sequences listing the vertices in this graph:
%%%  
%%% Random permutation: A, C, D, B
%%%
  \begin{compactenum}
% 1=A. BFS from vertex a %%% Verified 26-03-22
  \item%
    \label{sequence-bfs}%    
    $a$, $d$, $b$, $c$, $e$, $f$, $g$, $h$.
    
% 2=C. Dijkstra from vertex a %%% Verified 26-03-22
  \item%
    \label{sequence-dijkstra}%    
    $a$, $d$, $c$, $f$, $b$, $e$, $g$, $h$.
        
% 3=D. Spoiler        
  \item%
    \label{sequence-spoiler}%    
    $a$, $d$, $c$, $b$, $e$, $g$, $h$, $f$.
%%% Not Dijkstra, because b is further away than f   
%%% Not DFS, because after c a neighbour of c would have been visited   
%%% Not BFS, because we would have f after e, not g   
        
% 4=B. DFS from vertex a %%% Verified 26-03-22
  \item%
    \label{sequence-dfs}%
    $a$, $d$, $b$, $c$, $e$, $g$, $f$, $h$.

  \end{compactenum}
  %
  For each of the sequences above, determine whether it can be the result of:
  \begin{compactenum}[(a)]
  \item
    a breadth-first search with vertices listed in order of visits;
    
  \item
    a depth-first search with vertices listed in order of discovery;
    
  \item
    a shortest-path computation with vertices listed in the order they
    are removed from the priority queue.

  \end{compactenum}
%
  Each sequence above is the output of at most one of the algorithms,
  but since there are four sequences there could be a sequence that
  cannot be produced by any of the algorithms.

  
\begin{figure}[t] %%  [tp]
  \centering
  \includegraphics{Figures/graph-algorithm-outputs-26.1}
%
  \caption{Graph for 
    Problem~\ref{problem:graph-algorithm-outputs}%
  }
  \label{fig:graph-algorithm-outputs}
\end{figure}

    

  Partial credit is given for matching algorithms and vertex sequences correctly.
%
  For full credit, you need to
  also
  give a
  brief
  overview of how and why the
  algorithms process the vertices in the given order
  (such as explaining the order of recursive calls, edges processed,
  or similar),
  or why none of the proposed algorithms could yield the sequence in
  question,  
  but you do not have to provide detailed information about any
  auxiliary data structures used in the algorithms.

  \begin{solution}
    *** UPDATE ***
    
  Let us consider the different algorithms in the order listed and
  try to find sequences that would match their outputs.
  Since we are promised that each sequence is the output of at most
  one algorithm, we are done with each sequence as soon as we find a
  match.
  
  \textit{\textbf{Sequence~\ref{sequence-bfs}}}
  is the result of breadth-first search starting with vertex~$a$.
  To see this, consider how the different vertices are enqueued in and
  dequeued from the queue if we start the breadth-first search in
  vertex~$a$:

  \medskip
  \noindent
  \begin{tabular}[]{|l|l|l|}
    \hline
    \textbf{Dequeued vertex}
    &
      \textbf{Enqueued vertices}
    &
      \textbf{Queue}
    \\
    \hline
    ---
    &
      ---
    &
      $(a)$
    \\
    $a$
    &
      $\set{b, c, e}$
    &
      $(b, c, e)$
      \\
    $b$
    &
      $\set{d}$
    &
      $(c, e, d)$
    \\
    $c$
    &
      ---
    &
      $(e, d)$
    \\
    $e$
    &
      $\set{h}$
    &
      $(d, h)$
    \\
    $d$
    &
      ---
    &
      $(h)$
    \\
    $h$
    &
      $\set{g}$
    &
      $(g)$
    \\
    $g$
    &
      $\set{f}$
    &
      $(f)$
    \\
    $f$
    &
      ---
    &
      $()$
    \\
    \hline
  \end{tabular}

  \bigskip

  
  \textit{\textbf{Sequence~\ref{sequence-dfs}}}
  is the result of depth-first
  search starting with vertex~$a$.
  Below follows an illustration of in which orders vertices are
  discovered and finished: \\
  \vspace{-6.0mm}
  \begin{tabbing}
    Disco\={}ver  $a$ --- undiscovered neighbours $b$, $c$ and~$e$
    \\
    \> Disco\={}ver  $b$ --- undiscovered neighbour $d$
    \\
    \> \> Disco\={}ver  $d$  --- undiscovered neighbour $c$
    \\
    \> \> \> Disco\={}ver  $c$ --- no undiscovered neighbours
    \\
    \> \> \> Finish $c$, since no undiscovered neighbours
    \\
    \> \> Finish $d$, since no undiscovered neighbours left
    \\
    \> Finish $b$, since no undiscovered neighbours left
    \\
    \> Discover $e$ --- undiscovered neighbour $h$
    \\
    \> \> Discover $h$  --- undiscovered neighbour $g$
    \\
    \> \> \> Discover $g$  --- undiscovered neighbour $f$
    \\
    \> \> \> \> Discover $f$ --- no undiscovered neighbours
    \\
    \> \> \> \> Finish $f$, since no undiscovered neighbours 
    \\
    \> \> \> Finish $g$, since no undiscovered neighbours left
    \\
    \> \> Finish $h$, since no undiscovered neighbours left
    \\
    \> Finish $e$, since no undiscovered neighbours left
    \\
    Finish $a$, since no undiscovered neighbours left
  \end{tabbing}
  
  
  \textit{\textbf{Sequence~\ref{sequence-dijkstra}}}
  is the result of a call to Dijkstra's algorithm computing shortest
  paths from vertex~$a$.
  To see this, consider in which order the vertices are dequeued from
  the priority queue and how vertex key values are updated as edges
  are relaxed:

  \medskip
  \noindent
  \begin{tabular}[]{|l|l|l|}
    \hline
    \textbf{Dequeued}
    &
%         \textbf{Vertex updates}
      \textbf{Updates}
    &
      \textbf{Priority queue (after relaxations)}
    \\
    \hline
    ---
    &
      ---
    &
      $
      \set{
      (a: 0),
      (b: \infty),
      (c: \infty),
      (d: \infty),
      (e: \infty),
      (f: \infty),
      (g: \infty),
      (h: \infty)
      }
      $
    \\
    $a$
    &
      $\set{b, c, e}$
    &
      $
      \set{
      (c: 2),
      (e: 4),
      (b: 6),
      (d: \infty),
      (f: \infty),
      (g: \infty),
      (h: \infty)
      }
      $
    \\
    $c$
    &
      $\set{b}$
    &
      $
      \set{
      (e: 4),
      (b: 5),
      (d: \infty),
      (f: \infty),
      (g: \infty),
      (h: \infty)
      }
      $
    \\
    $e$
    &
      $\set{h}$
    &
      $
      \set{
      (b: 5),
      (h: 6),      
      (d: \infty),
      (f: \infty),
      (g: \infty)
      }
      $
    \\
    $b$
    &
      $\set{d}$
    &
      $
      \set{
      (h: 6),      
      (d: 12),
      (f: \infty),
      (g: \infty)
      }
      $
    \\
    $h$
    &
      $\set{g}$
    &
      $
      \set{
      (g: 7),      
      (d: 12),
      (f: \infty)
      }
      $
    \\
    $g$
    &
      $\set{f}$
    &
      $
      \set{
      (f: 9),      
      (d: 12)
      }
      $
    \\
    $f$
    &
      $\set{d}$
    &
      $
      \set{
      (d: 10)
      }
      $
    \\
    $d$
    &
      ---
    &
      $
      \set{
      }
      $
    \\
    \hline
  \end{tabular}

  \bigskip
  
  \textit{\textbf{Sequence~\ref{sequence-spoiler}}}, finally,
  is a spoiler sequence that cannot be the output of any of the
  algorithms listed. To see this, consider the following case
  analysis:
  \begin{compactitem}
  \item
    If the sequence were the result of breadth-first search,
    then it would list vertices in order of the number of directed
    edges required to get to the vertex     from
    the start vertex~$a$, but the first deviation from this is that
    $g$ is listed before~$d$.
    
  \item
    A depth-first search starting with $a$, $c$, $b$ would discover~$d$
    next and not~$e$.
        
  \item
    Finally, the sequence cannot be the output of Dijkstra's    algorithms,
    since the vertices are not listed in increasing order of distance
    from~$a$. The first deviation here is that $e$ should come
    before~$b$ in a listing in increasing order of distance.
    
  \end{compactitem}

  
\end{solution}

%     
\begin{figure}[t] %%  [tp]
  \centering
  \includegraphics{Figures/graph-algorithm-outputs-26.1}
%
  \caption{Graph for 
    Problem~\ref{problem:graph-algorithm-outputs}%
  }
  \label{fig:graph-algorithm-outputs}
\end{figure}

    

  
\end{problem}





\end{document}


